\section{The transition map for contact}\label{sec:tmap}

Every contact algorithm answers two questions about a query point $\xbf$ that sits near a body $X$. Is
$\xbf$ inside $X$, and if so by how much? This signed depth is the \emph{gap}. Which way is outward? This
direction is the \emph{normal}. The neural-atlas framework answers both questions from one of two ways of
describing the body. The thesis of this paper is that the description, not the scheme that enforces the
contact constraint, controls accuracy on geometrically hard contact.

A few conventions fix the signs and the notation we use throughout. The gap is negative when the point
penetrates the body: $\gn<0\Leftrightarrow$ penetration, so an interior point reads a negative gap and two
surfaces that overlap read a negative gap. The Macaulay bracket $\pos{z}=\max(0,z)$ keeps only the
positive part of its argument and so isolates the penetrating cases. We label bodies $A,B,\dots$; a body's
interior is the open region $\Omega_X\subset\R^d$ in $d=2$ or $3$ dimensions, and its boundary is the
surface $\partial\Omega_X$. We describe that boundary in two ways. A \emph{chart} parametrizes
the boundary directly, $\map{X}:\Theta_X\to\partial\Omega_X$, sending a parameter $\theta\in\Theta_X$ to
the surface point it labels. A \emph{level set} instead fills all of space with a scalar field
$\phi_X:\R^d\to\R$ whose zero crossing is the boundary $\partial\Omega_X$. The next subsection compares the
two.

\subsection{Two ways to describe a body}\label{sec:tmap:two}

\subsubsection{The ambient level set}

The classical description stores a signed-distance function, a scalar field over all of space whose value
is the distance to the surface, negative inside and positive outside. For body $A$ we train a neural
network $\phi_A:\R^d\to\R$ to approximate this field, and we read the gap and the outward normal from the
field and its gradient. The gap is the field value, and the normal is the field's gradient made into a
unit vector.
%
\begin{equation}\label{eq:sdf-gap}
\gn^A(\xbf)=\phi_A(\xbf),\qquad
\nbf_A(\xbf)=\frac{\grad\phi_A(\xbf)}{\norm{\grad\phi_A(\xbf)}}.
\end{equation}
%
Here $\grad\phi_A$ is the spatial gradient of the field and $\norm{\cdot}$ is the Euclidean norm, so both
the gap and the normal come from one automatic-differentiation call. We train the network to behave like a
true distance through the Eikonal penalty $\mathbb E_{\xbf}\big[(\norm{\grad\phi}-1)^2\big]$, which drives
the gradient to unit length, together with a surface loss, a normal-alignment loss, and sign-anchor terms
\citep{liu2020implicit,park2019deepsdf}; the architecture is a $\tanh$ multilayer perceptron.

The level set fails on hard geometry in two ways. Both are built into the description rather than being
fixable implementation defects. The first is the medial axis. The exact Euclidean signed distance is only
$C^0$ on the locus of points equidistant from two or more boundary points, because there the nearest-point
map is multivalued and the gradient $\grad\phi$ jumps. In a concavity this medial set reaches close to the
surface, so the gradient normal degrades exactly where a nonconvex contact happens
\citep{kolev2003level}. The second is spectral bias. A coordinate network learns low frequencies first and
high frequencies slowly or never, because its neural-tangent kernel is Laplace-like
\citep{rahaman2019spectral}. A network $\phi$ of fixed capacity therefore smooths cusps, asperities, and
fine detail; it cannot represent a boundary sharper than its capacity ceiling.
Section~\ref{sec:fourier} makes this quantitative.

\subsubsection{The boundary chart}

The alternative stores only the boundary, written as a function of the boundary's own parameter. Three
chart families appear in this work, in order of increasing generality. The \emph{radial chart}
$\rho_X:\Sph^{d-1}\to\R^+$ is a height field on the sphere of directions: for each unit direction it
records how far the surface lies from the center. It works for star-shaped bodies, those seen entirely
from one interior point. With center $c_X$ and orientation $Q_X$ (the body axes as columns), the boundary
point in direction $\hat\ubf$ is $\map{X}(\hat\ubf)=c_X+Q_X\,\rho_X(\hat\ubf)\,\hat\ubf$. The \emph{height
chart} $h_X:\mathcal D\to\R$ over an open parameter domain $\mathcal D$ is the graph $z=h(x)$ of a
single-valued rough surface, such as one face of a rock joint. The \emph{decoder chart}
$D_X:\bm\xi\mapsto\xbf$ is a boundary-fitted volumetric map for general, possibly non-star-shaped geometry,
inverted by Newton iteration.

A chart costs less to fit than a level set, for a simple reason. The radius or
height function lives in one dimension fewer than the ambient field, models no interior, and carries no
medial axis. The price is a precondition that the body must meet (star-shapedness for radial charts,
single-valuedness for height charts) and a gap measured along a ray rather than perpendicular to the
surface (Section~\ref{sec:tmap:bias}). The payoff is fidelity, because a $(d{-}1)$-dimensional boundary
function is far easier to fit sharply than a $d$-dimensional ambient field.
Figure~\ref{fig:tm-levelset-vs-chart} contrasts the two descriptions.

\begin{figure}
\centering
\resizebox{\linewidth}{!}{%
\begin{tikzpicture}[font=\footnotesize,>=Stealth,line join=round,
  starpts/.style={}]

  % ===== reusable star coordinates (concave notches) =====
  % defined relative to each panel origin via shifts

  %% ---------------------------------------------------------------
  %% PANEL (a): level set
  %% ---------------------------------------------------------------
  \begin{scope}[shift={(0,0)}]
    % nested OUTER level-set contours (fainter)
    \draw[gray!35,line width=0.4pt] plot[smooth cycle,tension=0.7]
      coordinates {(0:1.55)(45:0.95)(90:1.55)(135:0.95)(180:1.55)(225:0.95)(270:1.55)(315:0.95)};
    \draw[gray!50,line width=0.4pt] plot[smooth cycle,tension=0.7]
      coordinates {(0:1.4)(45:0.82)(90:1.4)(135:0.82)(180:1.4)(225:0.82)(270:1.4)(315:0.82)};
    % the body boundary (zero level set)
    \draw[black,line width=0.9pt] plot[smooth cycle,tension=0.7]
      coordinates {(0:1.25)(45:0.68)(90:1.25)(135:0.68)(180:1.25)(225:0.68)(270:1.25)(315:0.68)};
    % nested INNER contours (fainter)
    \draw[gray!50,line width=0.4pt] plot[smooth cycle,tension=0.7]
      coordinates {(0:1.05)(45:0.52)(90:1.05)(135:0.52)(180:1.05)(225:0.52)(270:1.05)(315:0.52)};
    \draw[gray!35,line width=0.4pt] plot[smooth cycle,tension=0.7]
      coordinates {(0:0.82)(45:0.38)(90:0.82)(135:0.38)(180:0.82)(225:0.38)(270:0.82)(315:0.38)};
    % medial axis: dashed line from a concave notch toward interior
    \draw[red!70,dashed,line width=0.7pt] (45:0.68) -- (0,0);
    \node[red!75,font=\scriptsize,anchor=west] at (0.18,0.62) {medial axis};
    % title
    \node[anchor=north,text width=3.2cm,align=center] at (0,-1.75)
      {(a) level set $\phi:\mathbb{R}^d\!\to\!\mathbb{R}$};
  \end{scope}

  %% ---------------------------------------------------------------
  %% PANEL (b): radial chart
  %% ---------------------------------------------------------------
  \begin{scope}[shift={(4.0,0)}]
    % body boundary in teal
    \draw[teal,line width=0.9pt] plot[smooth cycle,tension=0.7]
      coordinates {(0:1.25)(45:0.68)(90:1.25)(135:0.68)(180:1.25)(225:0.68)(270:1.25)(315:0.68)};
    % 8 thin rays center -> boundary
    \draw[teal!70,line width=0.4pt] (0,0)--(0:1.25);
    \draw[teal!70,line width=0.4pt] (0,0)--(45:0.68);
    \draw[teal!70,line width=0.4pt] (0,0)--(90:1.25);
    \draw[teal!70,line width=0.4pt] (0,0)--(135:0.68);
    \draw[teal!70,line width=0.4pt] (0,0)--(180:1.25);
    \draw[teal!70,line width=0.4pt] (0,0)--(225:0.68);
    \draw[teal!70,line width=0.4pt] (0,0)--(270:1.25);
    \draw[teal!70,line width=0.4pt] (0,0)--(315:0.68);
    % center dot
    \fill[teal] (0,0) circle (1.6pt);
    \node[teal!80,font=\scriptsize,anchor=north east] at (-0.04,-0.04) {$c$};
    % a sample angle label
    \node[teal!80,font=\scriptsize] at (0:1.0) [above] {$\rho(\theta)$};
    % title
    \node[anchor=north,text width=3.2cm,align=center] at (0,-1.75)
      {(b) radial chart $\rho:\mathbb{S}^1\!\to\!\mathbb{R}^{+}$};
  \end{scope}

  %% ---------------------------------------------------------------
  %% PANEL (c): the radial gap
  %% ---------------------------------------------------------------
  \begin{scope}[shift={(8.0,0)}]
    % body boundary
    \draw[teal,line width=0.9pt] plot[smooth cycle,tension=0.7]
      coordinates {(0:1.25)(45:0.68)(90:1.25)(135:0.68)(180:1.25)(225:0.68)(270:1.25)(315:0.68)};
    % center dot c
    \fill[black] (0,0) circle (1.6pt);
    \node[font=\scriptsize,anchor=north east] at (-0.04,-0.04) {$c$};
    % exterior point p along ~20 deg ray, beyond boundary
    \coordinate (P) at (20:1.95);
    % foot where ray meets boundary (green) ~ rho at 20deg
    \coordinate (F) at (20:1.15);
    % the ray c -> p
    \draw[gray!70,line width=0.5pt,dashed] (0,0)--(P);
    % unit direction marker hat d
    \node[font=\scriptsize,anchor=south] at (20:0.7) {$\hat{d}$};
    % foot dot (green)
    \fill[teal!60!black] (F) circle (1.8pt);
    % exterior point p
    \fill[orange] (P) circle (1.8pt);
    \node[font=\scriptsize,anchor=west] at ($(P)+(0.06,0.02)$) {$p$};
    % violet double arrow foot -> p  (the gap)
    \draw[violet,line width=0.9pt,<->] ($(F)+(20:0.02)$) -- ($(P)-(20:0.02)$);
    \node[violet,font=\scriptsize,anchor=south west] at ($(F)!0.5!(P)+(0,0.08)$)
      {$g=r-\rho(\hat{d})$};
    % title
    \node[anchor=north,text width=3.2cm,align=center] at (0,-1.75)
      {(c) the radial gap};
  \end{scope}

\end{tikzpicture}%
}
\caption{Two representations of the same star-shaped body. (a) The ambient level set
$\phi:\R^d\to\R$ fills all of space; its zero set is the boundary, but it carries a medial axis in the
interior where $\grad\phi$ is non-smooth, and a fixed-capacity neural $\phi$ low-passes sharp features.
(b) The radial boundary chart $\rho:\Sph^{1}\to\R^+$ stores only the boundary as a height field on the
circle of directions, with no ambient field and no medial axis. (c) The gap is read along the ray from
the center, $\gn=r-\rho(\dhat)$, with the matched normal $\grad_d F=\dhat-\grad_S\rho/r$.}
\label{fig:tm-levelset-vs-chart}
\end{figure}

\subsection{The transition map and why it detects contact}\label{sec:tmap:def}

\subsubsection{The idea before the formula}

Here is the central construction in one sentence. Take a point on body $A$'s surface and ask where that
point lands on body $B$'s surface. The transition map is the rule that answers, carrying every surface
point of $A$ to its counterpart on $B$. Two charts make this possible: each surface point of $A$ is named
by a parameter $\theta_A$, each surface point of $B$ by a parameter $\psi_B$, and the transition map turns
the first name into the second. It does so in two steps. First push the $A$-point into physical space with
$A$'s chart, then invert $B$'s chart to read off the $B$-parameter at that location.
%
\begin{equation}\label{eq:tmap-def}
\tmap:\ \theta_A\ \longmapsto\ \psi_B=\big(\map{B}^{-1}\circ\map{A}\big)(\theta_A).
\end{equation}
%
Reading the formula left to right: $\map{A}$ takes the parameter $\theta_A$ to the physical point
$\xbf=\map{A}(\theta_A)$ on $A$, and $\map{B}^{-1}$ takes that physical point back to the matching surface
parameter $\psi_B$ on $B$. The name is the standard one from differential geometry, where a transition map
is the change of coordinates $\map{B}^{-1}\circ\map{A}$ between two charts of an atlas; we apply it across
two bodies to read their contact correspondence. Figure~\ref{fig:tm-concept} illustrates the construction.

Once the correspondence is known, the contact normal follows from the chart's own geometry rather than
from a signed-distance gradient. Two tangent vectors span the surface at the contact point, and their
cross product points out of it. Normalizing that cross product gives the unit normal.
%
\begin{equation}\label{eq:jacobian-normal}
\nbf=\frac{\tbf_1\times\tbf_2}{\norm{\tbf_1\times\tbf_2}},\qquad
\tbf_\alpha=\frac{\partial\map{A}}{\partial\xi^\alpha},
\end{equation}
%
Here $\tbf_\alpha$ is the surface tangent obtained by differentiating the chart $\map{A}$ along its
parameter direction $\xi^\alpha$, and the normalized cross product $\nbf$ is the outward unit normal. The
undeformed normal gap is then the chart form of the implicit level-set MPM normal gap
\citep{liu2020implicit}, $g_{n0}=(\map{}^{(j)}-\map{}^{(i)})\cdot\nbf$, the projection of the
surface-to-surface offset onto that normal.

\begin{figure}
\centering
\resizebox{\linewidth}{!}{%
\begin{tikzpicture}[font=\footnotesize,>={Stealth[length=2.4mm]}]
  \fill[blue!12,draw=blue!70,line width=0.8pt]
    plot[smooth cycle,tension=0.8] coordinates
    {(-1.7,1.5) (-0.4,1.85) (0.95,1.0) (1.25,0.0) (0.95,-1.25)
     (-0.4,-1.7) (-1.7,-1.25) (-2.25,0.0) (-1.95,0.95)};
  \fill[orange!14,draw=orange!85!black,line width=0.8pt]
    plot[smooth cycle,tension=0.8] coordinates
    {(4.5,1.55) (5.85,1.9) (7.0,1.05) (7.35,0.0) (6.9,-1.3)
     (5.55,-1.7) (4.35,-1.2) (3.95,0.0) (4.15,1.0)};
  \coordinate (cA) at (-0.55,0.05);
  \coordinate (cB) at (5.55,0.0);
  \fill (cA) circle (1.3pt) node[below=1pt]{$c_A$};
  \fill (cB) circle (1.3pt) node[below=1pt]{$c_B$};
  \coordinate (x) at (1.1,0.2);
  \coordinate (foot) at (3.95,0.1);
  % rays
  \draw[blue!70,dashed,line width=0.7pt] (cA) -- (x) node[pos=0.45,above,sloped]{$\theta_A$};
  \draw[orange!85!black,dashed,line width=0.7pt] (cB) -- (foot) node[pos=0.55,below,sloped]{$\psi_B$};
  \fill (x) circle (1.6pt);
  \fill[orange!85!black] ($(foot)+(-0.07,-0.07)$) rectangle ($(foot)+(0.07,0.07)$);
  % gap arrow (long, between separated bodies)
  \draw[violet,line width=1pt,<->] (x) -- (foot);
  % spread labels: x-label up-left, gap formula high, foot-label down-right
  \node[anchor=east] at ($(x)+(-0.05,0.45)$) {$x=\varphi_A(\theta_A)$};
  \draw[gray,line width=0.3pt] ($(x)+(-0.05,0.40)$) -- (x);
  \node[violet,anchor=south] at (2.5,1.15) {$g_B(x)=\|x-c_B\|-\rho_B(\psi_B)$};
  \draw[violet,line width=0.3pt] (2.5,1.12) -- ($(x)!0.5!(foot)+(0,0.06)$);
  \node[anchor=north west] at ($(foot)+(0.02,-0.18)$) {$\varphi_B(\psi_B)$};
  % body labels + boxed map
  \node[font=\small\bfseries,blue!70!black] at (-0.55,-2.05) {body $A$};
  \node[font=\small\bfseries,orange!75!black] at (5.55,-2.05) {body $B$};
  \node[draw=black,line width=0.7pt,rounded corners=2pt,inner sep=5pt,font=\small]
    at (2.5,-2.95) {$\tau_{AB}:\ \theta_A\mapsto\psi_B=(\varphi_B^{-1}\circ\varphi_A)(\theta_A)$};
\end{tikzpicture}%
}
\caption{The transition map $\tmap=\map{B}^{-1}\circ\map{A}$. A boundary point $\xbf=\map{A}(\theta_A)$ of
body $A$ is taken into physical space, then $B$'s chart is inverted along the ray from $c_B$ to give
$\psi_B$ and the foot $\map{B}(\psi_B)$. The signed gap $\gn^B(\xbf)=\norm{\xbf-c_B}-\rho_B(\psi_B)$ falls
out of the inversion; no ambient field is evaluated.}
\label{fig:tm-concept}
\end{figure}

\subsubsection{Detection needs no ambient field}

% code: solvers/contact/chart_gap.py::evaluate_gap_chart ; solvers/contact/supershape.py::radial_gap
Detection asks only for a signed gap and an outward normal at a query point, and neither requires an
ambient field. Given $B$'s chart, three steps suffice: express the query point $\xbf$ in $B$'s local
coordinates, invert the chart to find the boundary parameter that the point points toward, and compare the
two radii. For a star-shaped body all three collapse to almost nothing: the inversion is a closed-form
direction, and the comparison is a single radius subtraction. We first form
the body-frame coordinate $\bm d=Q_B^\top(\xbf-c_B)$, which re-expresses the offset from $B$'s center in
$B$'s own axes, then split it into a direction and a length.
%
\begin{equation}\label{eq:radial-gap}
\dhat=\frac{Q_B^\top(\xbf-c_B)}{\norm{Q_B^\top(\xbf-c_B)}},\qquad
\gn^B(\xbf)=\underbrace{\norm{Q_B^\top(\xbf-c_B)}}_{r}-\rho_B(\dhat),
\end{equation}
%
Here $\dhat$ is the unit direction from $B$'s center to the query point, $r$ is the distance from the
center to the point, and $\rho_B(\dhat)$ is the chart's surface radius in that direction; the gap is the
point's radius minus the surface's radius. This gap is globally single-valued and $C^1$ wherever $\rho_B$
is, and it is negative exactly when the point lies inside. This is the level set turned inside out. A
signed-distance function answers ``how far is the nearest surface, in any direction'' by storing a field
over all of $\R^d$; the transition map answers ``which point of $B$'s boundary does this point correspond
to, and is it inside'' by storing only $\partial\Omega_B$ and inverting. On star-shaped and single-valued
geometry the second is both cheaper to store and sharper to fit.

In two dimensions the radial chart can carry an analytic radius, the Gielis superformula, which generates
a wide family of closed boundaries from a few parameters.
%
\begin{equation}\label{eq:superformula-chart}
\rho(\theta)=\Big(\big\lvert\tfrac1a\cos\tfrac{m\theta}{4}\big\rvert^{n_2}
+\big\lvert\tfrac1b\sin\tfrac{m\theta}{4}\big\rvert^{n_3}\Big)^{-1/n_1},
\end{equation}
%
Here $\theta$ is the polar angle, $m$ sets the number of lobes, and $a,b,n_1,n_2,n_3$ shape the lobes; the
inverse chart is simply $\theta=\angle\,Q_B^\top(\xbf-c_B)$, the angle of the body-frame offset, evaluated
by \code{atan2} \citep{gielis2003generic}. The analytic radius and a trained neural radius feed the
identical gap expression~\eqref{eq:radial-gap}: the evaluator itself is validated against the closed-form
gap to $10^{-10}$ on a self-test, and the trained neural radius reproduces the analytic gap to
$3.8\times10^{-3}L$ on the cusped superformula (Section~\ref{sec:verif}), so the neural chart is a drop-in
replacement for the analytic one.

\subsection{The matched conservative normal}\label{sec:tmap:matched}

One detail decides whether the contact forces conserve energy: the normal we return with the gap must be
the gradient of that gap. Only then is the penalty force the gradient of a potential, so that the work it
does is path-independent and nothing leaks. Differentiating the radial gap $F(\xbf)=r-\rho(\dhat)$ asks
for the gradient of $\rho$ with respect to direction, but that raw gradient is not what we want, because
$\rho$ is a function of a \emph{unit} vector. Moving along the direction $\dhat$ would change its length,
and $\rho$ does not respond to a length change. The raw gradient therefore carries a spurious
radial component that we must remove. We remove it by projecting onto the plane tangent to the direction
sphere.
%
\begin{equation}\label{eq:tangential-proj}
\grad_S\rho=(I-\dhat\dhat^\top)\,\grad_{\hat d}\rho,
\end{equation}
%
Here $\grad_{\hat d}\rho$ is the raw direction gradient of the radius, $I-\dhat\dhat^\top$ is the
projector that strips off the component along $\dhat$, and $\grad_S\rho$ is the surviving tangential
(surface) gradient. With the spurious part gone, the gap gradient in the body frame and the world-frame
normal both follow cleanly.
%
\begin{equation}\label{eq:matched-normal}
\grad_d F=\dhat-\frac{\grad_S\rho}{r},\qquad
\nbf=\widehat{\,Q\,\grad_d F\,}.
\end{equation}
%
Here $\grad_d F$ is the gap gradient in $B$'s body frame, the radial unit vector $\dhat$ corrected by the
tangential slope of the surface, and $\nbf$ is that vector rotated into world axes by $Q$ and normalized
(the hat denotes normalization). With this matched pair the penalty force is exactly
$-\grad\!\big(\tfrac12\epsilon_n\pos{-F}^2 V\big)$, the gradient of the penalty potential, and so is
conservative; an unmatched normal would leak energy. The projection~\eqref{eq:tangential-proj} is what
makes this work for any non-spherical body, where the slope $\grad_S\rho$ is nonzero.

The forces that consume the matched pair are standard. Penalty, augmented-Lagrangian, and
regularized-Coulomb friction all take only the triple $(\gn,\nbf,\norm{\fbf_N})$ as input
\citep{wriggers2006computational,kikuchi1988contact,laursen2002computational,alart1991mixed,simo1992augmented}.
Because they depend on nothing else, the chart and the level set are interchangeable at the force layer,
and swapping one description for the other changes only the geometry, never the constraint solver.

This interchangeability has a useful consequence on geometry both descriptions can represent. On a sphere
the radial chart reproduces the sphere signed distance, gap and normal alike, to machine precision at
all depths and orientations, and substituting the chart detector into a live two-sphere MPM collision is
bit-identical in velocity, impulse, and penetration. The chart never does worse than the level set where
both apply; it pulls ahead only where the level set degrades.

\subsection{Multi-arc enumeration}\label{sec:tmap:multiarc}

The transition map carries a kinematic advantage. Picture sliding a finger
along body $A$'s whole boundary and, at each step, asking whether that point has crossed into $B$. Wherever
the answer is yes, the finger is inside a contact region; the boundary breaks into runs of ``inside'' and
``outside,'' and each inside run is one contact patch. A single closest-point projection cannot see this
structure, because it returns only the one globally nearest foot. Concretely, we scan $A$'s chart
$\theta_A\mapsto\xbf=\map{A}(\theta_A)$, evaluate the inverse-chart gap~\eqref{eq:radial-gap} against $B$
at each $\theta_A$, and mark the parameters with $\gn^B<0$; the contiguous marked runs are the contact
arcs. For a star shape this is a one-dimensional sweep over $\theta$ in two dimensions, or a geodesic cap
scan in three dimensions.

Figure~\ref{fig:tm-multiarc} shows the payoff on a nonconvex pair. The boundary scan reports two or more
disjoint arcs where a single closest-point projection returns only the globally nearest foot and misses
the second contact entirely. This matters for cams, gears, and any geometry that touches in more than one
place. On the cusped superformula benchmark of Section~\ref{sec:verif} the chart scan reports $\ge 2$ arcs
where a single closest-point projection reports $1$. The obstacle chart must still be star-shaped;
conventional node-to-surface and mortar schemes recover disjoint patches by segment search instead
\citep{zavarise1998segment,puso2004mortar,fischer2005frictionless}.

\begin{figure}
\centering
\resizebox{0.9\linewidth}{!}{%
\begin{tikzpicture}[font=\footnotesize,>=Stealth,line join=round]
  % ---------- body B (gently wavy top, filled below) ----------
  \def\wavetop{
    (-4.2,0.45)
    .. controls (-3.3,0.95) and (-2.6,0.05) .. (-1.8,0.55)
    .. controls (-1.0,1.0)  and (-0.4,0.35) .. (0.3,0.6)
    .. controls (1.0,0.85)  and (1.7,0.15)  .. (2.4,0.5)
    .. controls (3.1,0.85)  and (3.6,0.25)  .. (4.2,0.55)
  }
  \fill[orange!18]
    \wavetop -- (4.2,-1.6) -- (-4.2,-1.6) -- cycle;
  \draw[orange!85!black,thick] \wavetop;
  \node[orange!60!black,font=\small\bfseries] at (3.3,-1.15) {body $B$};

  % ---------- follower A: two-footed staple ----------
  % circle centres
  \coordinate (cL) at (-1.2,1.35);
  \coordinate (cR) at ( 1.2,1.35);
  \def\rr{0.55}
  % top connecting bar
  \draw[blue!70!black,thick,fill=blue!10]
    (-1.2,1.95) -- (1.2,1.95) -- (1.2,2.25) -- (-1.2,2.25) -- cycle;
  % the two circular feet
  \fill[blue!10] (cL) circle (\rr);
  \fill[blue!10] (cR) circle (\rr);
  \draw[blue!70!black,thick] (cL) circle (\rr);
  \draw[blue!70!black,thick] (cR) circle (\rr);
  \node[blue!60!black,font=\small\bfseries] at (0,2.7) {follower $A$};

  % ---------- contact arcs dipping below the curve ----------
  % left foot: lower arc below curve (shallower contact)
  \draw[red,line width=1.4pt] ([shift=(232:\rr)]cL) arc (232:308:\rr);
  \node[red] at (-1.2,0.45) {arc 1};
  % right foot: lower arc below curve (deeper contact)
  \draw[red,line width=1.4pt] ([shift=(225:\rr)]cR) arc (225:315:\rr);
  \node[red] at (1.2,0.5) {arc 2};

  % ---------- single CPP on the deeper (right) contact ----------
  \coordinate (lp) at (1.2,0.8);              % lowest point of right foot
  \draw[teal!70!black,thick] (lp) circle (0.16);
  \node[teal!60!black,anchor=west] at (2.05,0.05)
    {single CPP: 1 foot};
  \draw[teal!70!black,->] (3.0,0.2) .. controls (2.2,0.4) .. (lp);
\end{tikzpicture}%
}
\caption{A boundary scan over body $A$'s chart evaluates the inverse-chart gap~\eqref{eq:radial-gap} at
every $\theta_A$ and marks the parameters with $\gn^B<0$. On a nonconvex pair these form two disjoint
contact arcs; a single closest-point projection returns only the globally nearest foot, missing the
second contact.}
\label{fig:tm-multiarc}
\end{figure}

\subsection{Inverting the chart in practice}\label{sec:tmap:invert}

% code: common/geometry.py::invert_decoder ; solvers/contact/profile_chart_2d.py::vertical_gap
The step that costs the most is inverting $B$'s chart, $\map{B}^{-1}$, and its cost depends on the chart
family. It takes three forms, in increasing order of work.

For \emph{radial charts} the inverse is closed form. The boundary parameter is just the direction angle
$\psi=\angle\,Q_B^\top(\xbf-c_B)$, computed by \code{atan2} in two dimensions or by normalizing
$Q_B^\top(\xbf-c_B)$ in three dimensions; the cost is $O(1)$ per query with no iteration.

For \emph{height charts} there is nothing to invert at all. The boundary is the graph $\map{}(x)=(x,h(x))$,
so the domain coordinate $x$ already \emph{is} the inverse-chart coordinate, and the transition map reads
off the gap by a single height comparison. The gap is how far the query point sits above the surface, and
the normal is perpendicular to the surface slope.
%
\begin{equation}\label{eq:height-gap}
g_{\mathrm{vert}}(x,z)=z-h_\theta(x),\qquad
\nbf=\frac{(-h_\theta'(x),\,1)}{\sqrt{1+h_\theta'^2(x)}}.
\end{equation}
%
Here $z$ is the height of the query point, $h_\theta(x)$ is the learned surface height at $x$, and
$h_\theta'(x)$ is its slope; the normal tilts away from vertical in proportion to that slope, and the
denominator normalizes it. The three-dimensional analogue is
$\nbf=(-h_x,-h_y,1)/\sqrt{1+\norm{\grad h}^2}$, built from the two surface slopes $h_x,h_y$. This is the
form the rough rock-joint detector of Section~\ref{sec:cv7} uses, and it is the reason that detector costs
$O(1)$ per query with no Newton solve.

For a general \emph{decoder chart} $D_B(\bm\xi)$ no closed form exists, so the inverse is found by Newton
iteration: we look for the parameter $\bm\xi$ whose image $D_B(\bm\xi)$ equals the query point $\xbf$, by
driving the residual $D_B(\bm\xi)-\xbf$ to zero.
%
\begin{equation}\label{eq:newton-invert}
\bm\xi^{(k+1)}=\bm\xi^{(k)}-J^{-1}\big(D_B(\bm\xi^{(k)})-\xbf\big),\qquad
J=\frac{\partial D_B}{\partial\bm\xi}.
\end{equation}
%
Here $\bm\xi^{(k)}$ is the $k$th guess, $J$ is the decoder Jacobian assembled by automatic
differentiation, and each step corrects the guess by the linearized residual. We stabilize the solve by
flooring the Jacobian's singular values and keeping the determinant consistent with the clamped spectrum,
which guards against a near-folded chart. The initial guess is the linear local-coordinate projection, and
the iteration runs until $\norm{D_B(\bm\xi)-\xbf}<10^{-8}$.

\subsection{A conservative-large but sign-exact gap}\label{sec:tmap:bias}

% code: solvers/contact/contact_manager.py::body_gap_normal
The radial inverse-chart gap buys its cheapness at one cost: it is not a true Euclidean distance. Because
it measures along the ray from the center rather than perpendicular to the surface, it overshoots the
perpendicular distance whenever the surface is tilted relative to the ray. The angle $\alpha$ is the tilt
between the ray direction $\dhat$ and the true surface normal at the radial foot, and $g_\perp$ is the
perpendicular (shortest) distance.
%
\begin{equation}\label{eq:radial-bias}
g_{\mathrm{rad}}=\frac{g_\perp}{\cos\alpha}+O(g^2)
\quad\Longrightarrow\quad
\abs{g_{\mathrm{rad}}}\ge\abs{g_\perp}\ \text{always}.
\end{equation}
%
Because $\cos\alpha\le1$, the radial gap $g_{\mathrm{rad}}$ is never smaller in magnitude than the
perpendicular gap $g_\perp$; the two agree only when the ray meets the surface head-on.
Figure~\ref{fig:tm-radial-bias} shows the geometry. The consequences are precise and benign. The sign of
$g_{\mathrm{rad}}$ is exact, so for star-shaped bodies the active set $\{\gn<0\}$ of penetrating points is
identical to the true level set's. The magnitude is conservative-large, so an exterior point reads a
slightly larger gap and a penetrating point a slightly deeper penetration; the penalty it produces is
marginally stiffer, never softer. The value equals the signed distance only on the sphere and in the
$g\to0$ limit.

The radial field is $C^1$ but not eikonal, meaning its gradient is continuous but not of unit length. In
deep concave valleys $\norm{\grad g}$ can reach $\sim10^4$, because the radial normal rotates rapidly there.
A perpendicular-foot refinement recovers the true normal-projected distance, but we use it only for
verification, because that foot jumps across the medial axis and would leak energy inside the integrator.

Swapping the chart detector in for the level set is a one-flag change. Declaring a body with the chart
detector routes its gap-normal dispatch to the radial evaluator instead of the signed-distance evaluator,
while the force laws, the broad-phase culler, and the MPM and FEM scatter all stay the same. Floors and
half-spaces, which are not star-shaped about any finite center, stay on the signed-distance path. The
transition map is therefore an opt-in detector that composes with the rest of the contact stack rather than
forking it. For height charts the matching statement is that~\eqref{eq:height-gap} is the small-slope-exact
normal-projected gap, not the closest-point distance; for those charts the fidelity story is the asperity
slopes, not the gap magnitude, as Section~\ref{sec:cv7} measures.

\begin{figure}
\centering
\resizebox{0.62\linewidth}{!}{%
\begin{tikzpicture}[font=\footnotesize,>=Latex,line join=round]
  % ---- geometry ----------------------------------------------------------
  % wall inclined ~20 deg from vertical: direction (sin20, cos20)
  \coordinate (W1) at (1.30,-2.20);   % wall bottom
  \coordinate (W2) at (3.45, 3.70);   % wall top
  % interior point c (lower-left)
  \coordinate (c)  at (-0.70,-0.60);
  % radial foot on the wall (ray crossing)
  \coordinate (F)  at (1.85,-0.50);
  % exterior point p along the ray c->F, beyond the wall
  \coordinate (p)  at (3.55, 0.62);

  % ---- shaded interior side of the wall ----------------------------------
  \begin{scope}
    \clip (-1.6,-2.4) rectangle (4.6,4.0);
    \fill[teal!12]
      ($(W1)!1!(W2)$) -- (W1) -- (-1.6,-2.4) -- (-1.6,4.0) --
      ($ (W2) + (-1.6,0) $) -- (W2) -- cycle;
  \end{scope}
  % cleaner: just fill the half-plane left of the wall
  \begin{scope}
    \clip (-1.6,-2.4) rectangle (4.6,4.0);
    \fill[teal!12] (W1) -- (W2) -- (-1.6,4.0) -- (-1.6,-2.4) -- cycle;
  \end{scope}

  % ---- the wall (body boundary) ------------------------------------------
  \draw[teal,line width=1.6pt] (W1) -- (W2)
        node[pos=0.93,right=2pt,teal] {wall $\partial\Omega$};

  % ---- outward unit normal n at the radial foot --------------------------
  % outward = to the right of the wall; normal dir (cos20, -sin20) rotated:
  \coordinate (N) at ($(F)+(0.94,-0.34)$);
  \draw[->,black,line width=0.9pt] (F) -- (N) node[above right=-1pt,black] {$n$};

  % ---- radial ray  c -> foot -> p ----------------------------------------
  \draw[violet,line width=1.0pt,->] (c) -- (p);
  % along-ray gap segment foot->p (drawn in green over the ray)
  \draw[teal!70!black] (F) -- (p);

  % ---- perpendicular drop from p to the wall -----------------------------
  % foot of perpendicular from p onto line (W1)-(W2)
  \coordinate (Q) at ($(W1)!(p)!(W2)$);
  \draw[red,line width=1.0pt] (p) -- (Q);
  \draw[red,fill=red] ($(Q)+(-0.06,-0.06)$) rectangle ($(Q)+(0.06,0.06)$);
  \node[red,below left=1pt and 0pt of Q,align=center] {perpendicular\\ foot};

  % ---- angle alpha between normal and ray --------------------------------
  \pic[draw=black,angle radius=0.95cm,angle eccentricity=1.28,
       "$\alpha$"black] {angle = N--F--p};

  % ---- points ------------------------------------------------------------
  \fill[black] (c) circle (2.2pt) node[below left=0pt,black] {$c$};
  \fill[black] (p) circle (2.2pt) node[above right=0pt,black] {$p$};
  \fill[teal!70!black] (F) circle (2.2pt);
  \node[teal!70!black,below=10pt,align=center] at ($(F)+(-0.15,-0.05)$)
       {radial foot\\ $\rho(\hat d)$};

  % ---- segment labels ----------------------------------------------------
  \node[teal!70!black] at ($(F)!0.5!(p)+(0.18,0.30)$) {$g_{\mathrm{rad}}$};
  \node[red] at ($(p)!0.5!(Q)+(0.34,0.02)$) {$g_\perp$};

  % ---- boxed identity ----------------------------------------------------
  \node[draw,rounded corners,inner sep=5pt,fill=gray!6,align=center,
        font=\small] at (1.2,-3.15)
       {$g_{\mathrm{rad}}=g_\perp/\cos\alpha
         \ \Rightarrow\ |g_{\mathrm{rad}}|\ge|g_\perp|$};
\end{tikzpicture}%
}
\caption{On an inclined flank the ray from the center meets the surface at angle $\alpha$ to the true
normal, so the radial foot and the perpendicular foot differ and
$g_{\mathrm{rad}}=g_\perp/\cos\alpha\ge g_\perp$. The sign is exact, which yields the active-set
equivalence with the level set; the magnitude is conservative-large.}
\label{fig:tm-radial-bias}
\end{figure}
